\documentclass[a4paper,11pt]{article}

% --- PAQUETES ---
\usepackage[utf8]{inputenc}
\usepackage[spanish]{babel}
\usepackage[T1]{fontenc}
\usepackage[default]{lato}
\usepackage[top=0.8cm, bottom=0.8cm, left=1.2cm, right=1.2cm]{geometry} % Márgenes ligeramente más compactos
\usepackage{xcolor}
\usepackage{hyperref}
\usepackage{fontawesome5}
\usepackage{enumitem}
\usepackage{titlesec}
\usepackage{parskip}

% --- COLORES ---
\definecolor{darkgrey}{RGB}{40,40,40}
\definecolor{accent}{RGB}{204, 36, 29} 

% --- CONFIGURACIÓN ---
\hypersetup{colorlinks=true, linkcolor=accent, urlcolor=accent}
\setlength{\parindent}{0pt}

% Espaciado de secciones más ajustado
\titlespacing*{\section}{0pt}{6pt}{2pt}
\titleformat{\section}{\large\bfseries\color{darkgrey}\uppercase}{}{0em}{}[\titlerule]

% --- COMANDOS PERSONALIZADOS (Espaciado reducido) ---
\newcommand{\entry}[4]{
  \vspace{1pt}
  {\bfseries #1} \textbullet\ \textit{#2} \hfill \textbf{\footnotesize #3} \\
  {\small #4} \par 
  \vspace{2pt}
}

\newcommand{\project}[3]{
  \vspace{1pt}
  {\bfseries #1} \hfill \textit{\footnotesize #2} \\
  {\small #3} \par
  \vspace{2pt}
}

\newcommand{\award}[3]{
  \vspace{1pt}
  {\bfseries #1} \hfill \textbf{\footnotesize #2} \\
  {\small #3} \par
  \vspace{1pt}
}

% --- DOCUMENTO ---
\begin{document}
\pagestyle{empty}

% --- ENCABEZADO ---
\begin{center}
    {\LARGE \textbf{Mario Enrique Ramírez Gallardo}} \\ \vspace{1pt}
    {\large Ingeniero de Software | Product Builder} \\ \vspace{3pt}
    
    \small
    \faEnvelope\ \href{mailto:ragm030119@gs.utm.mx}{ragm030119@gs.utm.mx} \quad \textbullet \quad
    \faPhone\ (+52) 951 119 6182 \quad \textbullet \quad
    \faGithub\ \href{https://github.com/p3p3p3k4z}{p3p3p3k4z} \quad \textbullet \quad
    \faLinkedin\ \href{https://www.linkedin.com/in/mario-ramirez-abab9a318/}{/in/mario-ramirez} \\
    \faGlobe\ \href{https://marioramirez-dev.vercel.app/}{marioramirez-dev.vercel.app}
\end{center}

\vspace{-8pt}

% --- SOBRE MÍ ---
\section*{Sobre mí}
\small Desarrollador de software con sólida base en Ingeniería en Computación y dominio de entornos Linux. Especializado en la arquitectura y construcción de aplicaciones Web y Móviles. Poseo experiencia práctica en el desarrollo de sistemas y la creación de componentes modulares. Destaco por mi capacidad analítica para resolver problemas complejos bajo presión, enfocándome en entregar código escalable, mantenible y soluciones orientadas a resultados.

% --- EXPERIENCIA ---
\section*{Experiencia Profesional}

\entry{Teteocan Technologies}{Pasante de Desarrollo de Software}{Abr 2025 -- Sept 2025}
{Arquitecté soluciones móviles multiplataforma con \textbf{Flutter} y \textbf{Dart}, construyendo componentes modulares que aceleraron el ciclo de desarrollo. Diseñé la infraestructura en \textbf{Firebase} con sincronización en tiempo real, estructurando las reglas de seguridad y optimizando el esquema de base de datos para reducir significativamente la latencia en consultas críticas, garantizando la escalabilidad de la app.}

% --- PROYECTOS ---
\section*{Proyectos Destacados}

\project{\href{https://github.com/p3p3p3k4z/a-u-ita}{Añi Ita (Soul of the Flower) \faGithub}}{Python, Web, Datos NASA}{
Motor de restauración ambiental con datos satelitales \textbf{Landsat/MODIS} para automatizar el diseño de "Islas Polinizadoras". Incluye modelos predictivos de flora y un "Modo Abeja" inmersivo en 3D para validación visual.}

\project{\href{https://github.com/p3p3p3k4z/huellazo-solana}{Huellazo (Solana LATAM Hackathon) \faGithub}}{React, Solana, Rust, NFTs}{
Plataforma de turismo sostenible que gamifica el impacto ecológico mediante blockchain. Desarrollé la integración mediante RPCs personalizados, obteniendo un **pase directo a la fase de incubación**.}

\project{\href{https://github.com/p3p3p3k4z/Haskell-ML-Cat}{Haskell ML Cat \faGithub}}{Haskell, Python, Hasktorch, IPC}{
Sistema de benchmarking de alto rendimiento. Arquitectura híbrida que conecta un backend funcional en Haskell con una interfaz en Python vía \textbf{IPC}, comparando modelos KNN frente a Deep Learning.}



% --- LOGROS ---
\section*{Logros y Premios}

\award{3er Lugar -- CLIHC (Conf. Latinoamericana de Interacción Humano-Computadora)}{2026}{
Reconocimiento internacional por innovación en diseño y desarrollo de interfaces centradas en el usuario e impacto social.}

\award{Finalista Top 20 -- HackODS UNAM}{2026}{
Desarrollo de soluciones tecnológicas y visualización de datos utilizando portales de datos abiertos (INEGI, UNAM) para los ODS.}

\award{Top 5 Global -- Bolt x MiniMax VIBE HACKERS}{2025}{
Desarrollo individual de "Polli Vision", visión artificial para identificación de especies en tiempo real en menos de 48h.}

\award{Ganador Local (4to Lugar) -- NASA Space Apps Challenge}{2024}{
Desarrollo de "Cozobi", solución para optimización agrícola basada en telemetría geoespacial.}

% --- HABILIDADES ---
\section*{Habilidades}
\small
\textbf{Lenguajes:} Python, JavaScript/TypeScript, Dart, SQL, C, Bash. \\
\textbf{Infraestructura y Sistemas:} Linux, Docker, Git, CI/CD. \\
\textbf{Web y Móvil:} React, Next.js, Vite, Tailwind CSS, Flutter. \\
\textbf{Intereses:} DevOps, Machine Learning, Web3, Desarrollo de Sistemas.

% --- EDUCACIÓN Y CERTIFICACIONES ---
\section*{Educación y Certificaciones}

\textbf{Ing. en Computación (Prevista 2026)} \hfill \textit{Universidad Tecnológica de la Mixteca} \par \vspace{3pt}

\textbf{Machine Learning Engineer \& Associate Python Developer} \hfill \textit{DataCamp} \par \vspace{3pt}

\textbf{Ruta Profesional DevOps} \hfill \textit{CódigoFacilito} \par \vspace{3pt}

\textbf{Especialización en Back-end} \hfill \textit{Oracle Next Education (ONE) G7} \par \vspace{3pt}

\textbf{OCI Foundations Associate 2025} \hfill \textit{Oracle} \par

\end{document}